\documentclass{article}
\usepackage[margin=1in]{geometry}
\usepackage{amsmath, amssymb, amsthm}
\usepackage{enumitem}

%Formatting and Spacing
\usepackage{setspace}
\setenumerate[1]{label = (\alph*), parsep = 1pt, topsep = 5pt}
\setenumerate[2]{label = \roman*.}
\setlength\parindent{0pt}
\linespread{1.10}

%Custom Commands
\newcommand{\hmwkTitle}{Homework 2}
\newcommand{\hmwkDueDate}{September 13, 2021}
\newcommand{\hmwkClass}{Honors Discrete Mathematics}
\newcommand{\hmwkClassInstructor}{Professor Gerandy Brito}
\newcommand{\hmwkAuthorName}{Sarthak Mohanty}

%Fancy Headers and Footers
\usepackage{fancyhdr}
\usepackage{extramarks}
\pagestyle{fancy}
\lhead{\hmwkAuthorName}
\chead{\hmwkClass \ (\hmwkClassInstructor)}
\rhead{\hmwkTitle}
\cfoot{\thepage}
\renewcommand\headrulewidth{0.4pt}
\renewcommand\footrulewidth{0.4pt}

\title{
    \vspace{2in}
    \textbf{\hmwkClass:\ \hmwkTitle}\\
    \normalsize\vspace{0.1in}\small{Due\ on\ \hmwkDueDate\ at 11:59pm}\\
    \vspace{0.1in}\large{\textit{\hmwkClassInstructor}}
    \vspace{3in}
    \author{\textbf{\hmwkAuthorName} \\ \\ No collaborators}
    \date{}
}

\begin{document}

\maketitle
\pagebreak

\subsection*{Question 1} 
    For the arguments below, label each of the statements and write the premises and the conclusion using propositional variables and connectives. Then, check the validity of each of the arguments. You should properly label the rule of inference you use on each step along with the lines to which you apply it.
    \begin{enumerate}
        \item ``Doug, a student of this class, knows how to write programs in JAVA. Everyone who knows how to write programs in JAVA can get a high-paying job. Therefore, someone in this class can get a high paying job."
        \item ``Somebody in this class enjoys whale watching. Every person who enjoys whale watching cares about ocean pollution. Therefore, there is a person in this class who cares about ocean pollution."
    \end{enumerate}

\subsection*{Solution}
    \begin{enumerate}
        \item Let $P(x)$ denote ``$x$ knows how to write programs in JAVA," and let $J(x)$ denote ``$x$ can get a high-paying job." Then the premises are $P(\text{Doug})$ and $(\forall x)(P(x) \Rightarrow J(x))$. The conclusion is $(\exists x)(J(x))$. The argument constructed below shows that our premises lead to the desired conclusion.
        \begin{center}
        \begin{tabular}{ll}
            \textbf{Step} & \textbf{Reason} \\
            1. $P(\text{Doug})$ & Premise \\
            2. $(\exists x)(P(x))$ & Existential Generalization using (1) \\
            3. $(\forall x)(P(x) \Rightarrow J(x))$ & Premise \\
            4. $(\exists x)(J(x))$ & Universal Modus Ponens using (2) and (3)
        \end{tabular}
        \end{center}
        
        \item Let $W(x)$ denote ``$x$ enjoys whale watching." Let $P(x)$ denote ``$x$ cares about ocean pollution." The premises are $(\exists x)(W(x))$ and $(\forall x)(W(x) \Rightarrow P(x))$. The conclusion is $(\exists x)(P(x))$. The argument below shows that our premises lead to the desired conclusion.
        \begin{center}
        \begin{tabular}{ll}
            \textbf{Step} & \textbf{Reason} \\
            1. $(\exists x)(W(x))$ & Premise \\
            2. $(\forall x)(W(x) \Rightarrow P(x))$ & Premise \\
            3. $(\exists x)(P(x))$ & Universal Modus Ponens using (1) and (2)
        \end{tabular}
        \end{center}
    \end{enumerate}

\subsection*{Question 2}
    Show that the following argument is valid.
    $$\begin{array}{ll}
        (1) & (p \land t) \Rightarrow (r \lor s) \\
        (2) & q \Rightarrow (u \land t) \\
        (3) & u \Rightarrow p \\
        (4) & \neg s \\
        \hline
        \therefore & q \Rightarrow r
    \end{array}$$

\subsection*{Solution}
    \begin{center}
    \begin{tabular}{ll}
        \textbf{Step} & \textbf{Reason} \\
        (1) $q \Rightarrow (u \land t)$ & Premise \\
        (2) $(u \land t) \Rightarrow u$ & Simplification \\
        %$q \Rightarrow u$ & Simplification using (4) \\
        (3) $u \Rightarrow p$ & Premise \\
        %$q \Rightarrow p$ & Hypothetical Syllogism using (5) and (6) \\
        %$q \Rightarrow t$ & Simplification using (4) \\
        (4) $q \Rightarrow (p \land t)$ & Hypothetical Syllogism using (1), (2), and (3) \\
        (5) $(p \land t) \Rightarrow (r \lor s)$ & Premise \\
        (6) $\neg s$ & Premise \\
        (7) $(p \land t) \Rightarrow r$ & Disjunctive Syllogism using (5) and (6) \\
        (8) $q \Rightarrow r$ & Hypothetical Syllogism using (4) and (7) \\
    \end{tabular}
    \end{center}

\subsection*{Question 3}
    Show that if $n$ is an integer and $3n^2 + 2$ is even, then $n$ is even.
    \begin{enumerate}
        \item using a proof by contradiction.
        \item using a proof by contrapositive.
    \end{enumerate}

\subsection*{Solution}
    \begin{enumerate}
        \item Suppose $n$ is an integer and $3n^2 + 2$ is even. We wish to show that $n$ is even, and shall do so using contradiction. Suppose $n$ is odd. Then we can pick an integer $k$ such that $n = 2k + 1$. Then $3n^2 + 2 = 3(2k + 2)^2 + 1 = 3(4k^2 + 8k + 4) + 1 = 12k^2 + 24k + 13 = 2(6k^2 + 12k + 6) + 1$. But then $3n^2 + 2$ is odd. This is a contradiction, therefore $n$ is even.
        \item The contrapositive of the statement is ``if $n$ is odd, then $n$ is not an integer or $3n^2 + 2$ is odd". Suppose $n$ is odd. Then we can pick an integer $k$ such that $n = 2k + 1$. Then $3n^2 + 2 = 3(2k + 2)^2 + 1 = 3(4k^2 + 8k + 4) + 1 = 12k^2 + 24k + 13 = 2(6k^2 + 12k + 6) + 1$, so $3n^2 + 2$ is odd. Hence the contrapositive statement is always true, so the original statement is true as well.
    \end{enumerate}

\subsection*{Question 4}
    Prove the \textbf{triangle inequality}, which states that if $x, y \in \mathbb{R}$ then $|x| + |y| \ge |x + y|$. Here, $|x|$ is the absolute value of $x$ defined as 
    $$|x| = 
    \begin{cases}
         x & \text{if } x \ge 0 \\
        -x & \text{if } x < 0
    \end{cases}$$
    \textit{(Hint: Do a proof by cases, considering all possible choices of $x$ and $y$ been positive, negative, or zero.)}

\subsection*{Solution}
    We start with the trivial fact $$a^2 + 2|a||b| + b^2 \ge a^2 + 2ab + b^2.$$ Now observe that for all $x \in \mathbb{R}$, we have $|x|^2 = x^2$. %Prove?
    Hence we can simplify our first equation to obtain $$(|a| + |b|)^2 \ge |a + b|^2.$$ Since both $|a| + |b|$ and $|a + b|$ are positive, we can square root both sides to obtain our desired equation: $$|x| + |y| \ge |x + y|.$$

\subsection*{Question 5}
    Recall that the set of irrational numbers is those elements in the difference $\mathbb{R} - \mathbb{Q}$.
    \begin{enumerate}
        \item Is the sum of two irrational numbers always irrational? Either prove this statement or give a counterexample.
        \item Let $r \in \mathbb{Q}$ and $s \in \mathbb{R} - \mathbb{Q}$ be any two numbers, one rational and another irrational. Prove that $r + s$ is irrational. You may use the fact that the sum of two rational numbers is rational.
        \item Let $a \in \mathbb{Q}$ and $r \in \mathbb{R} - \mathbb{Q}$ be any two numbers, one rational and another irrational. You are also told that $a < r$. Show that there exists a number $x$ such that $a < x < r$ and $x$ is irrational. \textit{(Hint: there is a constructive proof for this problem. Part (b) may be helpful.)}
    \end{enumerate}

\subsection*{Solution}
    \begin{enumerate}
        \item No, $\pi$ and $-\pi$ are both irrational but their sum, $\pi + (-\pi) = 0$, is not.
        \item Suppose $r + s$ is rational. Since $r$ is rational, $-r$ must also be rational. Then since the sum of two rational numbers is rational, $(r + s) + (-r) = s$ is also rational. But $s$ is irrational. This is a contradiction, therefore $r + s$ is irrational.
        \item From part (b), we know that $a + r$ is irrational. Now $a < \frac{a + r}{2} < r$, so if $\frac{a + r}{2}$ is irrational, we are done.
        
        \vspace{1.5mm}
        Suppose $\frac{a + r}{2}$ is rational. Then there exists $b, c \in \mathbb{Z}$ such that $c \ne 0$ and $\frac{a + r}{2} = \frac{b}{c}$, so $a + r = \frac{2b}{c}$ is rational as well. But $a + r$ is irrational. This is a contradiction, therefore $\frac{a + r}{2}$ is irrational.
    \end{enumerate}

\end{document}